\chapter{Livrables du projet}

\section{Résultats attendus}

À l'issue du projet, les résultats concrets suivants seront livrés et opérationnels :
\begin{itemize}
    \item Une application Web React accessible sur le réseau intranet de l'ANTIC avec éditeur stylistique Excel ;
    \item Un moteur de traitement backend léger Python FastAPI robuste et évolutif ;
    \item Une base de données PostgreSQL ou MySQL configurée pour la journalisation d'audit ;
    \item Une suite de conteneurs Docker facilitant le déploiement et la maintenance de la solution ;
    \item Un document complet d'homologation et de recette fonctionnelle.
\end{itemize}

\section{Liste des livrables du projet}

\renewcommand{\arraystretch}{1.3}
\begin{table}[H]
\centering
\small
\begin{tabular}{|L{3.2cm}|L{4.2cm}|L{6.1cm}|}
\hline
\rowcolor{primarycolor!15}
\textbf{Domaine} & \textbf{Livrables Majeurs} & \textbf{Description \& Contenu} \\
\hline
\textbf{Spécifications} & Cahier des Charges \& Dossier d'Architecture & Spécifications fonctionnelles, modèle conceptuel de données (MCD) et architecture N-Tiers. \\
\hline
\textbf{Logiciel} & Application Web React \& Backend FastAPI & Interface composant React avec éditeur graphique Excel et API Python FastAPI. \\
\hline
\textbf{Infrastructure} & Conteneurs Docker \& Docker Compose & Scripts \texttt{Dockerfile} et \texttt{docker-compose.yml} préconfigurés pour PostgreSQL/MySQL et API. \\
\hline
\textbf{Documentation \& Support} & Manuel Utilisateur \& Administration & Guide illustré pour l'utilisation des agents et la gestion système de l'ANTIC. \\
\hline
\end{tabular}
\caption{Tableau synthétique des livrables du projet}
\label{tab:livrables_projet}
\end{table}
